\documentclass[10pt,a4paper]{amsart}
\usepackage[margin=19mm]{geometry}
\usepackage{amsmath,amssymb,mathtools}
\usepackage[scr=rsfs,cal=euler]{mathalfa}
\usepackage{xfrac}
\usepackage[unicode,hidelinks,bookmarks=false]{hyperref}
\newcommand{\ie}{i.e.\ }
\newcommand{\cf}{cf.\ }
\newcommand{\resp}{respectively\ }
\newcommand{\Subsection}{\subsection}
\providecommand{\nobreakdash}{\nobreak\hbox{-}}
\setlength{\emergencystretch}{2em}
\multlinegap0pt
\usepackage{amssymb}
\usepackage{natbib}
\usepackage{multirow}
\usepackage{bm}
\usepackage{xcolor}
\usepackage{cancel}
\usepackage{mathtools}
\usepackage{tikz}
\usepackage{caption}
\usepackage{subfig}
\newtheorem{theorem}{Theorem}
\title{The true detection probability versus the subjective detection probability of a uniformly optimal search plan}
\author{Liang Hong}
\date{Source: arXiv:2601.00350v1 (CC BY 4.0)}
\begin{document}
\maketitle

\noindent\emph{Source and licence.} The true detection probability versus the subjective detection probability of a uniformly optimal search plan. Version arXiv:2601.00350v1 (CC BY 4.0), distributed by arXiv under the Creative Commons Attribution 4.0 International licence, http://creativecommons.org/licenses/by/4.0/, as recorded in the arXiv licence metadata for this exact version and shown on the abstract page https://arxiv.org/abs/2601.00350v1. Full licence notice: \texttt{licenses/arxiv-2601.00350-CC-BY-4.0.txt}.\par\smallskip

\noindent\textit{Editorial scope.} Counted unit: the article's theorem thm:1discrete (source0.tex lines 290--293). The article's complete original proof of it is delivered with it (source0.tex lines 295--361, one \texttt{proof} environment of 67 lines). No mathematics is written, repaired or re-derived: the statement and the proof are the article's own lines, in the article's own notation, and no retained line is substituted or re-flowed.  No further source-local context is needed: every cross-reference of every retained line resolves inside the two delivered spans.  Nothing was omitted from any retained span, so the package carries no omission record.  No retained line contains a citation, so the wrapper carries no bibliography and no bibliography entry.  No retained line contains a figure environment, an image-inclusion command or drawing code, so no image is delivered and the package declares no asset dependency.  Editorial formatting only, in addition: an AMS \texttt{amsart} preamble that restates the article's own declarations and macros used by the retained lines, namely the environments \texttt{theorem} on separate counters, together with the standard packages those lines need.  The retained environments print in this document as Theorem 1 in order of appearance, whereas the article prints the same items with the numbering of its own full document.  The complete native source of the article as distributed in the arXiv e-print for this exact version is bundled unchanged as \texttt{sources/191993/source0.tex}.
\begin{theorem}
\label{thm:1discrete}
Suppose the possibility area $X$ is discrete, $x_0\in X$,  and the target distribution is uniform on $X$.  If the detection function $d$ is regular and homogeneous, then $P[\varphi^\star(\cdot, t)]=P^{\#}[\varphi^\star(\cdot, t)]$ for all $t\geq 0$ and $\mu(\varphi^\star)=\mu^{\#}(\varphi^\star)$.
\end{theorem}

\begin{proof}
Without loss of generality, we assume that $X=\{1, 2, \ldots, m\}$.  By assumption, the target density function is given by 
\[
\pi(x)=\left\{
		                           \begin{array}{ll}
		                           \frac{1}{m},& \hbox{for $x\in X$,} \\
					0,  & \hbox{otherwise.} 
		                          \end{array}
		                         \right.
\]
Since $d$ is homogeneous,  we have $\frac{\partial}{\partial y}d(x,y)=d'(y)$ and 
\[
q_x(y)=\left\{
		                           \begin{array}{ll}
		                           \frac{1}{m} d'(y),& \hbox{if $x\in D$,} \\
					0,  & \hbox{otherwise,} 
		                          \end{array}
		                         \right.
\]
Therefore,
\begin{eqnarray*}
q_x^{-1}(\lambda) &=& \left\{
		                           \begin{array}{ll}
		                          \text{inverse of $d'(y)$ evaluated at $\frac{\lambda}{\pi(x)}$},& \hbox{for $0<\lambda \leq q_x(0)$ and $x\in D$,} \\
					0,  & \hbox{otherwise,} 
		                          \end{array}
		                         \right.\\
		                         &=& \left\{
		                           \begin{array}{ll}
		                          \text{inverse of $d'(y)$ evaluated at $\lambda m$},& \hbox{for $0<\lambda \leq \frac{d'(0)}{m}$ and $x\in D$,} \\
					0,  & \hbox{otherwise,} 
		                          \end{array}
		                         \right.
\end{eqnarray*}
and 
\[
Q(\lambda)=\sum_{x\in X} q_x^{-1}(\lambda)=m q_x^{-1}(\lambda).
\]
Therefore, 
\[
Q^{-1}(K)=\left(q_x^{-1}\right)^{-1}\left(K/m\right)=q_x \left(K/m \right),
\]
where the amount of effort $K\geq 0$.  It follows that
\begin{eqnarray*}
\varphi^\star(x, t)&=&
q_x^{-1}(Q^{-1}(E(t)))\\
&=&\left\{
		                           \begin{array}{ll}
		                           q_x^{-1}\left(q_x \left(E(t)/m \right)\right),& \hbox{for $0<q_x \left(E(t)/m\right) \leq q_x(0)$ and $x\in D$,} \\
					0,  & \hbox{otherwise,} 
		                          \end{array}
		                         \right.\\
&=&\left\{
		                           \begin{array}{ll}
		                         \frac{E(t)}{m}, & \hbox{$x\in D$,} \\
					0,  & \hbox{otherwise,} 
		                          \end{array}
		                         \right.		                         
\end{eqnarray*}
where the last equality follows from the fact that $0<q_x \left(E(t)/V_D\right) \leq q_x(0)$ always holds (because $q_x(y)=\pi(x)d'(y)$ is strictly decreasing in $y$ for each $x\in X$). Thus, 
\begin{eqnarray*}
P[\varphi^\star(\cdot, t)] &=& \sum_{x\in X} d(x, \varphi^\star(x, t))\pi(x)
                                         = \sum_{x \in X} d(\varphi^\star(x, t))\pi(x) =\sum_{x\in X} d(E(t)/m)\pi(x) \\
                                         &= &d(E(t)/m)=d(x_0, \varphi^\star(\cdot, t))=P^{\#}[\varphi^\star(\cdot, t)].
\end{eqnarray*}
Hence,  $\mu(\varphi^\star)=\mu^{\#}(\varphi^\star)$.
\end{proof}

\end{document}
